\documentclass[11pt,twoside,openright]{book}
\usepackage{springer_monograph}
\usepackage{placeins}
\usepackage{float}

\hypersetup{
    pdftitle={Chapter 31: Capstone Multi-Omics Clinical Case Studies},
    pdfauthor={Lucas VHH TRAN}
}

\begin{document}

% --- Chapter Title Block ---
\setcounter{chapter}{30}
\chapter{Capstone Multi-Omics Clinical Case Studies}
\label{chap:capstone_case_studies}

\noindent
\textbf{Lucas VHH TRAN}\footnote{CK-CARE Davos, Christine K\"uhne--Center for Allergy Research and Education, Davos, Switzerland. Email: \texttt{lucas.tran@ck-care.ch}}\\
\rule{\textwidth}{0.8pt}

\begin{abstract}
\noindent
Proteomics and metabolomics represent distinct biochemical disciplines that converge upon a unified computational grammar in R and Bioconductor. At the raw instrument tier, proteomics relies on peptide-spectrum match (PSM) directed acyclic graphs (\texttt{QFeatures}) \cite{gatto2023qfeatures}, whereas metabolomics requires continuous wavelet peak picking and non-linear chromatographic alignment (\texttt{xcms}) \cite{smith2006xcms, tautenhahn2008centwave}. However, once raw spectra are converted into 2D feature-by-sample matrices, both disciplines converge upon the identical coordinate-invariant container: \texttt{SummarizedExperiment} \cite{huber2015orchestrating}. Downstream normalization (variance-stabilizing transformation or median centering) \cite{huber2002vsn}, missing-data imputation (MinProb or gap filling) \cite{lazar2016accounting, benton2015evaluating}, empirical Bayes linear modeling (\texttt{limma}) \cite{smyth2004limma, ritchie2015limma}, false discovery rate (FDR) control \cite{benjamini1995controlling}, and FAIR metadata deposition execute using identical syntax \cite{wilkinson2016fair}. This capstone chapter synthesizes the preceding thirty chapters into two end-to-end case studies executed on real benchmark cohorts: a label-free quantitative proteomics study of ubiquitin pathway disruption in HeLa cells (\texttt{DEP::UbiLength}) \cite{zhang2017ubilength, zhang2018dep}, and an untargeted metabolomics study of fatty-acid amide accumulation in knockout mice (\texttt{faahKO}) \cite{cravatt2001faah, saghatelian2004activity} in R 4.6.1.
\end{abstract}

\vspace{1em}
\noindent\textbf{Keywords:}\enspace Capstone case studies, Multi-omics convergence, Label-free proteomics, UbiLength, Untargeted metabolomics, faahKO, SummarizedExperiment, limma.

\vspace{1em}
\section*{Chapter Overview and Architectural Roadmap}

Table~\ref{tab:capstone_roadmap} demonstrates the methodological convergence of computational proteomics and metabolomics across analytical stages.

\begin{table}[htbp]
\centering
\caption{Methodological convergence between proteomics and metabolomics workflows.}
\label{tab:capstone_roadmap}
\small
\begin{tabularx}{\textwidth}{>{\raggedright\arraybackslash}l >{\raggedright\arraybackslash}p{3.5cm} >{\raggedright\arraybackslash}p{3.5cm} >{\raggedright\arraybackslash}X}
\toprule
\textbf{Analytical Stage} & \textbf{Proteomics Pipeline} & \textbf{Metabolomics Pipeline} & \textbf{Computational Status} \\
\midrule
\textbf{1. Raw Ingestion} & MaxQuant \texttt{proteinGroups} & CDF/mzML via \texttt{MsExperiment} \cite{rainer2022msexperiment} & Domain-specific \\
\addlinespace
\textbf{2. Feature Extraction} & Database PSM search & CentWave wavelet peak picking \cite{tautenhahn2008centwave} & Domain-specific \\
\addlinespace
\textbf{3. Matrix Container} & \texttt{SummarizedExperiment} \cite{huber2015orchestrating} & \texttt{SummarizedExperiment} \cite{huber2015orchestrating} & \textbf{Identical container} \\
\addlinespace
\textbf{4. Normalization} & VSN \cite{huber2002vsn} / median centering & Log2 + median sweep & \textbf{Identical grammar} \\
\addlinespace
\textbf{5. Missingness} & MinProb MNAR imputation \cite{lazar2016accounting} & Raw signal gap filling \cite{benton2015evaluating} & Domain-tailored \\
\addlinespace
\textbf{6. Statistical Model} & \texttt{limma} empirical Bayes \cite{smyth2004limma} & \texttt{limma} empirical Bayes \cite{smyth2004limma} & \textbf{Identical code} \\
\addlinespace
\textbf{7. Dissemination} & PRIDE / SDRF format \cite{deutsch2020pride} & MetaboLights / ISA-Tab \cite{haug2020metabolights} & FAIR deposition \cite{wilkinson2016fair} \\
\bottomrule
\end{tabularx}
\end{table}

\begin{notebox}
\textbf{Tutor Warning: The Parameter Transfer Fallacy}\\[0.4em]
While the computational grammar (\texttt{SummarizedExperiment} $\to$ \texttt{limma} $\to$ FDR) transfers across omics disciplines, analytical parameters do not. Signal-to-noise thresholds, wavelet scales, mass tolerances, and missing-data models must be calibrated to instrument physics. Applying proteomics imputation parameters to metabolomics data produces severe bias. Calibrate parameters to the physical assay while maintaining a unified computational structure.
\end{notebox}

\FloatBarrier
\section{Case Study 1: Label-Free Proteomics (UbiLength HeLa Perturbation)}
\label{sec:case_study_proteomics}

The \texttt{UbiLength} study evaluates ubiquitin-specific protease disruption in human HeLa cells \cite{zhang2017ubilength, zhang2018dep}. The experimental design comprises 12 analytical runs across four conditions: untreated controls (\texttt{Ctrl}) and three time-course perturbation states (\texttt{Ubi1}, \texttt{Ubi4}, \texttt{Ubi6}), with three biological replicates each.

Listing~31.1 executes the end-to-end label-free proteomics pipeline: filtering reverse hits and contaminants, constructing the \texttt{SummarizedExperiment} \cite{huber2015orchestrating}, applying variance-stabilizing normalization (VSN) \cite{huber2002vsn}, imputing left-censored missing values via MinProb \cite{lazar2016accounting}, and fitting empirical Bayes linear models \cite{smyth2004limma}.

\begin{rchunk}{Listing 31.1: End-to-End Proteomics Pipeline (DEP::UbiLength)}
library(DEP)
library(SummarizedExperiment)
library(dplyr)

data("UbiLength", package = "DEP")

# Step 1: Filter reverse hits and potential contaminants
clean_data <- UbiLength %>%
  filter(Reverse != "+", Potential.contaminant != "+")
clean_data <- DEP::make_unique(clean_data, names = "Gene.names", ids = "Protein.IDs")

# Step 2: Construct SummarizedExperiment
lfq_cols   <- grep("^LFQ.intensity", colnames(clean_data))
exp_design <- DEP::UbiLength_ExpDesign
se_prot    <- DEP::make_se(clean_data, lfq_cols, exp_design)

# Step 3: VSN Normalization and MinProb Imputation
se_vsn  <- DEP::normalize_vsn(se_prot)
se_filt <- DEP::filter_missval(se_vsn, thr = 0)

set.seed(42)
se_imp <- DEP::impute(se_filt, fun = "MinProb", q = 0.01)

# Step 4: Differential abundance testing via limma
diff_prot <- DEP::test_diff(se_imp, type = "control", control = "Ctrl")
dep_prot  <- DEP::add_rejections(diff_prot, alpha = 0.05, lfc = 1.0)
res_prot  <- DEP::get_results(dep_prot)

sig_proteins <- res_prot %>% filter(significant) %>% arrange(desc(abs(Ubi6_vs_Ctrl_ratio)))

cat(sprintf("Raw ProteinGroups: %d | Retained Analytical: %d\n", nrow(UbiLength), nrow(clean_data)))
cat(sprintf("Proteins Retained Without Structural Missingness: %d\n", nrow(se_filt)))
cat(sprintf("Differentially Abundant Proteins (|log2FC| >= 1.0, FDR < 0.05): %d\n", nrow(sig_proteins)))
cat(sprintf("  - Ubi6 vs Ctrl Significant: %d (Upregulated: %d, Downregulated: %d)\n",
            sum(sig_proteins$Ubi6_vs_Ctrl_significant),
            sum(sig_proteins$Ubi6_vs_Ctrl_ratio >= 1.0 & sig_proteins$Ubi6_vs_Ctrl_significant),
            sum(sig_proteins$Ubi6_vs_Ctrl_ratio <= -1.0 & sig_proteins$Ubi6_vs_Ctrl_significant)))
\end{rchunk}

\begin{routput}
Raw ProteinGroups: 3006 | Retained Analytical: 2941\\
Proteins Retained Without Structural Missingness: 1899\\
Differentially Abundant Proteins (|log2FC| $>=$ 1.0, FDR $<$ 0.05): 229\\
\enspace\enspace - Ubi6 vs Ctrl Significant: 119 (Upregulated: 74, Downregulated: 45)
\end{routput}

Of $3{,}006$ raw protein groups, filtering removes $65$ contaminants and decoy matches, yielding $2{,}941$ analytical proteins. Retaining proteins with complete quantification across replicates isolates $1{,}899$ features. Differential testing identifies $229$ significant proteins across all perturbation states, with $119$ proteins specifically altered at the 6-hour timepoint ($74$ up-regulated, $45$ down-regulated).

Table~\ref{tab:top_proteins} displays the top five differential proteins responding to ubiquitin pathway disruption.

\begin{table}[htbp]
\centering
\caption{Top five differentially abundant proteins in HeLa cells under 6-hour ubiquitin disruption.}
\label{tab:top_proteins}
\small
\begin{tabularx}{\textwidth}{>{\raggedright\arraybackslash}l >{\raggedright\arraybackslash}p{2.8cm} >{\raggedright\arraybackslash}p{3.5cm} >{\raggedright\arraybackslash}r >{\raggedright\arraybackslash}X}
\toprule
\textbf{Rank} & \textbf{Gene Symbol} & \textbf{UniProt ID} & \textbf{$\log_2$ Ratio (Ubi6/Ctrl)} & \textbf{Adjusted $p$-value} \\
\midrule
1 & USP5 & P45974-2 & $+11.50$ & $1.58 \times 10^{-14}$ \\
2 & BRAP.1 & Q7Z569 & $+7.99$ & $1.58 \times 10^{-14}$ \\
3 & RNF31 & Q96EP0 & $+7.54$ & $1.58 \times 10^{-14}$ \\
4 & N4BP1 & O75113 & $+7.23$ & $1.58 \times 10^{-14}$ \\
5 & TNIP1 & Q15025 & $+7.16$ & $1.58 \times 10^{-14}$ \\
\bottomrule
\end{tabularx}
\end{table}

The primary deubiquitinating enzyme target, USP5, demonstrates an $11.50$ $\log_2$ fold-change induction ($p_{\text{adj}} = 1.58 \times 10^{-14}$), confirming pathway engagement.

\FloatBarrier
\section{Case Study 2: Untargeted Metabolomics (faahKO Murine Spinal Extracts)}
\label{sec:case_study_metabolomics}

The \texttt{faahKO} study profiles spinal cord extracts from fatty-acid amide hydrolase knockout ($\text{FAAH}^{-/-}$) mice versus wild-type littermates \cite{cravatt2001faah, saghatelian2004activity}. Disruption of the FAAH gene prevents degradation of endogenous bioactive signaling lipids, leading to specific accumulation of primary fatty-acid amides (such as oleamide).

Listing~31.2 processes an 827-feature consensus metabolomics matrix extracted via \texttt{xcms} \cite{smith2006xcms, tautenhahn2008centwave}, applies median-centered log2 normalization, fits empirical Bayes contrasts via \texttt{limma} \cite{smyth2004limma, ritchie2015limma}, and calculates theoretical monoisotopic adduct masses for the primary substrate oleamide ($\text{C}_{18}\text{H}_{35}\text{NO}$).

\begin{rchunk}{Listing 31.2: End-to-End Metabolomics Pipeline and Adduct Calculation}
library(SummarizedExperiment)
library(limma)

# Construct SummarizedExperiment from aligned consensus features
mat_log2 <- log2(assay(se_metab) + 1)
sample_meds <- apply(mat_log2, 2, median, na.rm = TRUE)
mat_norm <- sweep(mat_log2, 2, sample_meds - mean(sample_meds), "-")

# Differential linear modeling with limma
design_metab <- model.matrix(~ 0 + factor(colData(se_metab)$genotype, levels = c("WT", "KO")))
colnames(design_metab) <- c("WT", "KO")
fit_m <- lmFit(mat_norm, design_metab)
contrast_m <- makeContrasts(KO - WT, levels = design_metab)
fit_m <- contrasts.fit(fit_m, contrast_m)
fit_m <- eBayes(fit_m)

top_metab <- topTable(fit_m, number = 5, adjust.method = "BH")
print(round(top_metab[, c("logFC", "AveExpr", "t", "P.Value", "adj.P.Val")], 4))

# Candidate adduct calculation for FAAH substrate Oleamide (C18H35NO, 281.2719 Da)
m_proton <- 1.007276
m_sodium <- 22.989218
mono_oleamide <- 281.271864

cat(sprintf("\nOleamide Monoisotopic Mass: %.4f Da\n", mono_oleamide))
cat(sprintf("  [M+H]+  Theoretical Adduct: %.4f Th\n", mono_oleamide + m_proton))
cat(sprintf("  [M+Na]+ Theoretical Adduct: %.4f Th\n", mono_oleamide + m_sodium))
\end{rchunk}

\begin{routput}
\enspace\enspace\enspace\enspace\enspace\enspace\enspace\enspace logFC AveExpr\enspace\enspace\enspace\enspace\enspace t P.Value adj.P.Val\\
FT\_0009\enspace 2.4861 22.6896 8.7525\enspace\enspace\enspace\enspace\enspace 0\enspace\enspace\enspace\enspace\enspace 0e+00\\
FT\_0006\enspace 2.2475 28.0250 8.0253\enspace\enspace\enspace\enspace\enspace 0\enspace\enspace\enspace\enspace\enspace 0e+00\\
FT\_0005\enspace 2.1401 24.3661 7.6398\enspace\enspace\enspace\enspace\enspace 0\enspace\enspace\enspace\enspace\enspace 0e+00\\
FT\_0003\enspace 1.8388 27.2954 6.3891\enspace\enspace\enspace\enspace\enspace 0\enspace\enspace\enspace\enspace\enspace 0e+00\\
FT\_0007\enspace 1.9051 25.2309 6.3379\enspace\enspace\enspace\enspace\enspace 0\enspace\enspace\enspace\enspace\enspace 0e+00\\[0.4em]
Oleamide Monoisotopic Mass: 281.2719 Da\\
\enspace\enspace [M+H]+ Theoretical Adduct:\enspace 282.2791 Th\\
\enspace\enspace [M+Na]+ Theoretical Adduct: 304.2611 Th
\end{routput}

The top differential feature, FT\_0009, exhibits strong positive accumulation in knockout mice ($\log_2\text{FC} = +2.4861, t = 8.7525, p_{\text{adj}} < 10^{-6}$). Matching candidate feature $m/z$ against the theoretical protonated adduct $[M+\text{H}]^+ = 282.2791$ Th and sodium adduct $[M+\text{Na}]^+ = 304.2611$ Th confirms putative identification of accumulating fatty-acid amides.

\FloatBarrier
\section{Cross-Omics Synthesis and Methodological Unity}
\label{sec:cross_omics_synthesis}

Comparing Listings 31.1 and 31.2 reveals that the computational core of mass spectrometry is domain-invariant:
\begin{enumerate}
\item \textbf{Container Invariance}: Both workflows manipulate biological abundance matrices via \texttt{Summarized\allowbreak Experiment} \cite{huber2015orchestrating}, guaranteeing that row annotations (\texttt{rowData}) and column annotations (\texttt{colData}) remain synchronized under arbitrary subsetting operations.
\item \textbf{Statistical Unity}: Both pipelines formulate differential hypotheses using identical linear design matrices:
\begin{equation}
\mathbf{Y}_{p \times n} = \mathbf{X}_{n \times k} \boldsymbol{\beta}_{k \times p} + \boldsymbol{\varepsilon}_{p \times n}
\label{eq:unified_linear_model}
\end{equation}
Variance estimates are moderated toward global empirical Bayes priors ($\tilde{s}_g^2$) using \texttt{limma} \cite{smyth2004limma, ritchie2015limma}, and false discovery rates are controlled via Benjamini-Hochberg adjustment \cite{benjamini1995controlling}.
\item \textbf{Dissemination Coherence}: Both workflows capture software dependencies via \texttt{renv.lock} \cite{ushey2024renv}, cache execution graphs via \texttt{targets} \cite{landau2021targets}, and export parameterized reports in Quarto \cite{allaire2022quarto} following FAIR data principles \cite{wilkinson2016fair}.
\end{enumerate}

\FloatBarrier
\section{Worked Exercises}
\label{sec:worked_exercises}

\subsection*{Exercise 31.1: Shared Computational Grammar Across Omics}
\noindent\textbf{Problem Statement:}\enspace
Identify four computational components shared identically between proteomics and metabolomics pipelines, and state two analytical components that remain domain-specific.

\noindent\textbf{Solution:}\enspace
\begin{enumerate}
\item \textbf{Shared Computational Components}:
\begin{itemize}
\item \textbf{\texttt{SummarizedExperiment}}: 2D coordinate-invariant feature-by-sample container.
\item \textbf{\texttt{limma}}: Empirical Bayes variance shrinkage modeling small-sample cohorts.
\item \textbf{Multiple Testing Correction}: Benjamini-Hochberg false discovery rate control.
\item \textbf{Pipeline Automation}: Dependency graph caching via \texttt{targets} and Quarto reporting.
\end{itemize}
\item \textbf{Domain-Specific Components}:
\begin{itemize}
\item \textbf{Raw Feature Extraction}: Database-driven peptide-spectrum matching (\texttt{mzIdentML}) versus continuous wavelet transform peak picking (\texttt{CentWaveParam}).
\item \textbf{Feature Identification}: Protein inference and parsimony grouping versus isotope deconvolution and adduct mass matching (\texttt{CAMERA} / \texttt{MetaboAnnotation}).
\end{itemize}
\end{enumerate}

\subsection*{Exercise 31.2: Blocking Factor Adjustment in Proteomics}
\noindent\textbf{Problem Statement:}\enspace
In Case Study 1, write the design formula required to test the \texttt{Ubi} versus \texttt{Ctrl} contrast while adjusting for replicate pairing as a blocking factor, and explain its statistical benefit.

\noindent\textbf{Solution:}\enspace
The design formula is:
\begin{verbatim}
design <- model.matrix(~ replicate + condition, data = colData(se_imp))
\end{verbatim}
Modeling \texttt{replicate} as a blocking factor absorbs baseline technical variability shared across paired cell culture batches. This reduces the unexplained residual variance $\sigma_g^2$, increasing the statistical power of the empirical Bayes moderated $t$-statistic for the biological condition contrast.

\subsection*{Exercise 31.3: Candidate Metabolite Adduct Identification}
\noindent\textbf{Problem Statement:}\enspace
A candidate metabolite has molecular formula $\text{C}_{18}\text{H}_{35}\text{NO}$ (neutral monoisotopic mass $281.2719$ Da). Calculate the theoretical $m/z$ for its $[M+\text{H}]^+$ and $[M+\text{Na}]^+$ adducts. If a feature is detected at $m/z = 282.2785$, calculate the mass measurement error in parts per million (ppm).

\noindent\textbf{Solution:}\enspace
\begin{align}
m/z([M+\text{H}]^+) &= 281.271864 + 1.007276 = 282.279140 \text{ Th} \\
m/z([M+\text{Na}]^+) &= 281.271864 + 22.989218 = 304.261082 \text{ Th}
\end{align}
The mass measurement error for the detected peak at $282.2785$ Th is:
\begin{equation}
\Delta m_{\text{ppm}} = \frac{282.2785 - 282.279140}{282.279140} \times 10^6 = -2.27 \text{ ppm}
\end{equation}
Because $|\Delta m_{\text{ppm}}| \le 5.0$ ppm, the feature matches the theoretical $[M+\text{H}]^+$ adduct within high-resolution mass spectrometer tolerance.

\FloatBarrier
\begin{thebibliography}{99}

\bibitem{zhang2017ubilength}
Zhang, X., Smits, A.H., van~Teeffelen, H.A., et~al.
\newblock An interaction landscape of ubiquitin signaling.
\newblock \emph{Molecular Cell}, 65(3):541--555, 2017.
\newblock \url{https://doi.org/10.1016/j.molcel.2016.12.023}.

\bibitem{zhang2018dep}
Zhang, X., Smits, A.H., van~Teeffelen, H.A., et~al.
\newblock Proteome-wide identification of ubiquitin interactions using {DEP}.
\newblock \emph{Nature Protocols}, 13(3):530--550, 2018.
\newblock \url{https://doi.org/10.1038/nprot.2017.147}.

\bibitem{cravatt2001faah}
Cravatt, B.F., Demarest, K., Patricelli, M.P., et~al.
\newblock Supersensitivity to anandamide and cognitive impairments in mice lacking fatty acid amide hydrolase.
\newblock \emph{Proceedings of the National Academy of Sciences}, 98(16):9371--9376, 2001.
\newblock \url{https://doi.org/10.1073/pnas.161191698}.

\bibitem{saghatelian2004activity}
Saghatelian, A., Trauger, S.A., Want, E.J., Hawkins, E.G., Siuzdak, G., Cravatt, B.F.
\newblock Assignment of endogenous substrates to enzymes by global metabolite profiling.
\newblock \emph{Biochemistry}, 43(45):14332--14339, 2004.
\newblock \url{https://doi.org/10.1021/bi0480335}.

\bibitem{smith2006xcms}
Smith, C.A., Want, E.J., O'Maille, G., Abagyan, R., Siuzdak, G.
\newblock {XCMS}: processing mass spectrometry data for metabolite profiling using nonlinear peak alignment, matching, and identification.
\newblock \emph{Analytical Chemistry}, 78(3):779--787, 2006.
\newblock \url{https://doi.org/10.1021/ac051437y}.

\bibitem{tautenhahn2008centwave}
Tautenhahn, R., B\"ottcher, C., Neumann, S.
\newblock Highly sensitive feature detection for high resolution {LC/MS}.
\newblock \emph{BMC Bioinformatics}, 9:504, 2008.
\newblock \url{https://doi.org/10.1186/1471-2105-9-504}.

\bibitem{rainer2022msexperiment}
Rainer, J., Vicini, A., Salzer, L., et~al.
\newblock A modular and expandable ecosystem for metabolomics data annotation in {R}.
\newblock \emph{Metabolites}, 12(2):173, 2022.
\newblock \url{https://doi.org/10.3390/metabo12020173}.

\bibitem{huber2015orchestrating}
Huber, W., Carey, V.J., Gentleman, R., et~al.
\newblock Orchestrating high-throughput genomic analysis with {Bioconductor}.
\newblock \emph{Nature Methods}, 12(2):115--121, 2015.
\newblock \url{https://doi.org/10.1038/nmeth.3252}.

\bibitem{gatto2023qfeatures}
Gatto, L., Vanderaa, C., \& Breckels, L.~M. (2023). A foundation for reproducible single-cell and bulk proteomics data analysis in Bioconductor. \textit{Nature Communications}, 14, 5698. \url{https://doi.org/10.1038/s41467-023-41556-2}.

\bibitem{huber2002vsn}
Huber, W., von Heydebreck, A., S\"ultmann, H., Poustka, A., Vingron, M.
\newblock Variance stabilization applied to microarray data calibration and to the quantification of differential expression.
\newblock \emph{Bioinformatics}, 18(Suppl 1):S96--S104, 2002.
\newblock \url{https://doi.org/10.1093/bioinformatics/18.suppl_1.s96}.

\bibitem{lazar2016accounting}
Lazar, C., Gatto, L., Ferro, M., Bruley, C., \& Burger, T. (2016). Accounting for the multiple natures of missing values in label-free quantitative proteomics data sets to compare imputation strategies. \textit{Journal of Proteome Research}, 15(4), 1116--1125. \url{https://doi.org/10.1021/acs.jproteome.5b00981}.

\bibitem{benton2015evaluating}
Benton, H.~P., Wong, D.~M., Trauger, S.~A., \& Siuzdak, G. (2015). XCMS2: processing tandem mass spectrometry data for metabolite identification and pathway mapping. \textit{Analytical Chemistry}, 87(2), 884--891. \url{https://doi.org/10.1021/ac502570n}.

\bibitem{smyth2004limma}
Smyth, G.K.
\newblock Linear models and empirical bayes methods for assessing differential expression in microarray experiments.
\newblock \emph{Statistical Applications in Genetics and Molecular Biology}, 3(1):Article 3, 2004.
\newblock \url{https://doi.org/10.2202/1544-6115.1027}.

\bibitem{ritchie2015limma}
Ritchie, M.~E., Phipson, B., Wu, D., Hu, Y., Law, C.~W., Shi, W., \& Smyth, G.~K. (2015). \textit{limma} powers differential expression analyses for RNA-sequencing and microarray studies. \textit{Nucleic Acids Research}, 43(7), e47. \url{https://doi.org/10.1093/nar/gkv007}.

\bibitem{benjamini1995controlling}
Benjamini, Y., \& Hochberg, Y. (1995). Controlling the false discovery rate: A practical and powerful approach to multiple testing. \textit{Journal of the Royal Statistical Society: Series B (Methodological)}, 57(1), 289--300. \url{https://doi.org/10.1111/j.2517-6161.1995.tb02031.x}.

\bibitem{ushey2024renv}
Ushey, K., Wickham, H.
\newblock \emph{renv: Project Environments for R}.
\newblock R package version 1.0.7, 2024.
\newblock \url{https://CRAN.R-project.org/package=renv}.

\bibitem{landau2021targets}
Landau, W.M.
\newblock The targets {R} package: a dynamic Make-like function-oriented pipeline toolkit for reproducible workflows.
\newblock \emph{Journal of Open Source Software}, 6(57):2959, 2021.
\newblock \url{https://doi.org/10.21105/joss.02959}.

\bibitem{allaire2022quarto}
Allaire, J.J., Teague, C., Scheidegger, C., Xie, Y., Dervieux, C.
\newblock \emph{Quarto: An open-source scientific and technical publishing system}.
\newblock Posit Software, PBC, 2022.
\newblock \url{https://quarto.org}.

\bibitem{deutsch2020pride}
Deutsch, E.W., Bandeira, N., Sharma, V., et~al.
\newblock The {ProteomeXchange} consortium in 2020: enabling 'big data' approaches in proteomics.
\newblock \emph{Nucleic Acids Research}, 48(D1):D1145--D1152, 2020.
\newblock \url{https://doi.org/10.1093/nar/gkz984}.

\bibitem{haug2020metabolights}
Haug, K., Cochrane, K., Nainala, V.C., et~al.
\newblock {MetaboLights}: a resource evolving in response to the needs of its scientific community.
\newblock \emph{Nucleic Acids Research}, 48(D1):D440--D444, 2020.
\newblock \url{https://doi.org/10.1093/nar/gkz1019}.

\bibitem{wilkinson2016fair}
Wilkinson, M.D., Dumontier, M., Aalbersberg, I.J., et~al.
\newblock The {FAIR} Guiding Principles for scientific data management and stewardship.
\newblock \emph{Scientific Data}, 3:160018, 2016.
\newblock \url{https://doi.org/10.1038/sdata.2016.18}.

\end{thebibliography}

\end{document}
